\section{Fazit}
	Wir kennen nun die Bedeutung der Joukowski-Transformation und können grafisch die Zirkulation bestimmen, welche für ein beliebiges Joukowski-Profil entsteht. Mit dieser Zirkulation kann unter Berücksichtigung der Flügellänge, der Luftdichte sowie der Strömungsgeschwindigkeit auf den Auftrieb eines Flügels geschlossen werden. Wichtig anzumerken ist hierbei, dass der Luftwiderstand zu Null resultiert. Somit werden die Luftreibung an den Flugzeugoberflächen sowie die Verwirbelungen vernachlässigt. Des Weiteren hat die Joukowski-Transformation eine weitere Einschränkung: Die Flügelhinterkante besitzt einen Winkel von 0 Grad, was logischerweise physikalisch nicht realisierbar ist. Obschon Joukowski-Flügelprofile nicht direkt in der Aviatik eingesetzt werden, stellt die Joukowski-Transformation einen der wichtigsten Meilensteine dar. Es war erstmals möglich, den Auftrieb einer Tragflächenform exakt zu berechnen.